\documentclass[10pt,a4paper]{amsart}
\usepackage[margin=19mm]{geometry}
\usepackage{amsmath,amssymb,mathtools}
\usepackage[scr=rsfs,cal=euler]{mathalfa}
\usepackage{xfrac}
\usepackage[unicode,hidelinks,bookmarks=false]{hyperref}
\newcommand{\ie}{i.e.\ }
\newcommand{\cf}{cf.\ }
\newcommand{\resp}{respectively\ }
\newcommand{\Subsection}{\subsection}
\providecommand{\nobreakdash}{\nobreak\hbox{-}}
\setlength{\emergencystretch}{2em}
\multlinegap0pt
\usepackage{amssymb}
\usepackage{xcolor}
\usepackage{tikz}
\usepackage{float}
\usepackage[shortlabels]{enumitem}
\newtheorem{lemma}{Lemma}
\newcommand{\ipr}[2]{\left\langle#1,#2\right\rangle}
\title{Sampling recovery in $L_2$ and other norms}
\author{David Krieg, Kateryna Pozharska, Mario Ullrich, Tino Ullrich}
\date{Source: arXiv:2305.07539v5 (CC BY 4.0)}
\begin{document}
\maketitle

\noindent\emph{Source and licence.} Sampling recovery in $L_2$ and other norms. Version arXiv:2305.07539v5 (CC BY 4.0), distributed by arXiv under the Creative Commons Attribution 4.0 International licence, http://creativecommons.org/licenses/by/4.0/, as recorded in the arXiv licence metadata for this exact version and shown on the abstract page https://arxiv.org/abs/2305.07539v5. Full licence notice: \texttt{licenses/arxiv-2305.07539-CC-BY-4.0.txt}.\par\smallskip

\noindent\textit{Editorial scope.} Counted unit: the article's lemma lem:lift (source0.tex lines 921--931). The article's complete original proof of it is delivered with it (source0.tex lines 934--997, one \texttt{proof} environment of 64 lines). No mathematics is written, repaired or re-derived: the statement and the proof are the article's own lines, in the article's own notation, and no retained line is substituted or re-flowed.  Retained uncounted as the source-local context the proof needs: lines 474--478; lines 735--739.  Nothing was omitted from any retained span, so the package carries no omission record.  No retained line contains a citation, so the wrapper carries no bibliography and no bibliography entry.  No retained line contains a figure environment, an image-inclusion command or drawing code, so no image is delivered and the package declares no asset dependency.  Editorial formatting only, in addition: an AMS \texttt{amsart} preamble that restates the article's own declarations and macros used by the retained lines, namely the environments \texttt{lemma} on separate counters, together with the standard packages those lines need.  The retained environments print in this document as Lemma 1 in order of appearance, whereas the article prints the same items with the numbering of its own full document.  The complete native source of the article as distributed in the arXiv e-print for this exact version is bundled unchanged as \texttt{sources/141424/source0.tex}.
\begin{equation}\label{eq:an}
\alpha_n \,:=\, \alpha(V_n,F, L_2
) \,:=\, \sup_{f \in F} \Vert f - P_{V_n} f  \Vert_2
\;=\; \sup_{f \in F}\, \inf_{g\in V_n}\Vert f - g  \Vert_2, 
\end{equation}

\begin{equation}\label{eq:Lambda}
\Lambda_n \;:=\;
\Lambda(V_n,G) \;:=\; 
 \sup\limits_{f \in V_n, \, f\neq 0} \frac{\|f\|_G}{\|f\|_2}.  
\end{equation}

\begin{lemma}\label{lem:lift}
 Let \hyperlink{assum}{Assumption~B} hold and let 
 $(\alpha_k)_{k=1}^\infty$ and $(\Lambda_k)_{k=1}^\infty$ be defined as in~\eqref{eq:an} and \eqref{eq:Lambda}, respectively.
 For any $m\in\mathbb{N}$, any mapping $A\colon F \to V_m$ and all $f\in F$,  we have
 \begin{equation*}%\label{eq:liftF}
  \Vert f - Af\Vert_{G}\,\le\, 
  2 \sum_{k > \lfloor m/4 \rfloor 
    } \frac{ \alpha_k \, \Lambda_{4k}}{k} 
  \,+\, \Lambda_m \cdot \Vert f-Af\Vert_2.
 \end{equation*}
 \end{lemma}

 \begin{proof}
Note that there is nothing to prove if the right hand side of the inequality % \eqref{eq:liftF} 
is infinite
 so we may assume that it is finite.
We first observe
\begin{align*}
 \Vert f - A f  \Vert_{G}
&\le\, \Vert f - P_{V_m} f \Vert_{G} + \Vert P_{V_m} f - A f \Vert_{G}\\
&\le\, \Vert f - P_{V_m} f \Vert_{G} + \Lambda_m \cdot \Vert P_{V_m}f - A f \Vert_{2}\\
&\le\, \Vert f - P_{V_m} f \Vert_{G} + \Lambda_m \cdot \Vert f - A f \Vert_{2}.
%{L_{2}} .
\end{align*}


In order to estimate $\Vert f - P_{V_m} f \Vert_G$,
we first 
show that $P_{V_m} f$ converges in $G\cap L_2$.
We 
%denote
 {define}
$I_\ell = \{ k \in \mathbb{N} \colon 2^{\ell}<  k \leq 2^{\ell+1}\}$ for $\ell\in\mathbb{N}_0$. Let $m\in I_s$ and $n\in I_t$ with $s\le t$. Then

\begin{align*}
\Vert P_{V_n} &f - P_{V_m} f \Vert_G  \,\le\,
\sum_{\ell=s}^t
\bigg\Vert\sum\limits_{ k\in I_\ell\colon m<k\le n} \ipr{f}{b_k}_2 b_k \bigg\Vert_G
\\ 
&\le\,
\sum_{\ell=s}^t
\bigg\Vert\sum\limits_{ k\in I_\ell} \ipr{f}{b_k}_2 b_k \bigg\Vert_2
  \cdot \Lambda_{2^{\ell+1}}
 \,\leq\, \sum_{\ell \geq s}  \alpha_{2^\ell} \Lambda_{2^{\ell+1}}.
\end{align*}
 Further we use the fact that the sequence $(\alpha_k)_{k=1}^\infty$
% {n\in \mathbb{N}}$ 
 is non-increasing, and therefore
\begin{equation*}%\label{eq:bound}
\alpha_{2^\ell} \leq   \frac{1}{2^{\ell-1}} \sum\limits_{j\in I_{\ell-1}}  {\alpha_j,}
	\end{equation*} 
and hence, since the sequence $(\Lambda_k)_{k=1}^\infty$
is non-decreasing,
\[
\begin{split}
\Vert P_{V_n} f - &P_{V_m} f \Vert_G
 \,\leq\, 2\sum_{\ell \geq s}  \sum\limits_{ j\in I_{\ell-1}}
\frac{\alpha_j \Lambda_{2^{\ell+1}}}{2^{\ell} } 
\\
& \leq 2\sum_{\ell \geq s}  \sum\limits_{ j\in I_{\ell-1}}
\frac{\alpha_j \Lambda_{4j}}{j}   
\,\le\, 2 \sum_{j > \lfloor m/4 \rfloor} \frac{ \alpha_j \Lambda_{4j}}{j},
\end{split}
\]
which is convergent due to the assumption.
This means that $(P_{V_m} f)$ is a Cauchy sequence in $G$.
Since we also know that $P_{V_m} f$ converges to $f$ in $L_2$, 
we obtain that  $(P_{V_m} f)$ is a Cauchy sequence in $G\cap L_2$ and converges to some $g\in G\cap L_2$ by Assumption (B.3). In particular, $f=g$ almost everywhere.
 Thus, using again Assumption (B.3), we obtain 
\[
 \Vert f - P_{V_m} f \Vert_G \,=\,
 \Vert g - P_{V_m} f \Vert_G \,\le\, 
 2 \sum_{j > \lfloor m/4 \rfloor} \frac{ \alpha_j \Lambda_{4j}}{j},
\]
which yields the statement.
 \end{proof}

\end{document}
